\section{Measure the exact output before selecting a capacity}
Let $H(u)$ count the bytes in the shortest supported head for unsigned argument $u$.
It is one through nine according to the width table. For a signed reading $x$, use
argument $A(x)=x$ when $x\geq0$ and $A(x)=-1-x$ otherwise. With counter $T$, identifier
length $L$ and $N$ readings, the exact output size is
\[
 4+H(T)+H(L)+L+H(N)+\sum_{i=0}^{N-1}H(A(x_i)).
\]
The four fixed bytes are the root map head and its three one-byte keys. Identifier
contents do not affect this size, but sample values do. The bound entry therefore
reads the sixteen-byte prefix and the selected sample codes; it need not read any
identifier bytes.

For the trace, the terms are four, nine, one, three, one and nineteen, giving thirty-seven.
An empty record with counter zero occupies seven. The largest deterministic result in
this profile occupies 241 bytes: nine counter-head bytes, thirty-two identifier bytes,
two-byte length heads for identifier and array, and sixty-four three-byte readings,
plus the four fixed bytes. A capacity of thirty-seven must accept the worked record;
requiring the maximum 241 for every record would reject a sufficient buffer.
